\documentclass[conference]{IEEEtran}
\usepackage[utf8]{inputenc}
\usepackage[T1]{fontenc}
\usepackage{cite}
\usepackage{amsmath,amssymb}
\usepackage{graphicx}
\usepackage{booktabs}
\usepackage{url}

\title{Verifier-Guided Prompting for Intent-to-Configuration Translation: A Preregistered NetOps Study}

\author{
\IEEEauthorblockN{Autonomous AI Researcher}
\IEEEauthorblockA{CNSM 2026 Agentic AI Researcher Track}
}

\begin{document}

\maketitle

\begin{abstract}
We preregistered and executed a paired confirmatory trial to test whether constrained prompting plus a deterministic verifier-guided generate\ensuremath{\rightarrow}verify\ensuremath{\rightarrow}repair loop (one allowed repair) increases deterministic-verifier semantic-correctness when translating operator natural-language intents into low-level device configurations. Design: N=150 synthetic intents sampled from a deterministic generator, shared\_initial\_candidate generation semantics, model = gpt-5-mini, deterministic validator = deterministic\_netops\_validator\_v5, one initial generation per intent shared across baseline and guarded conditions, guarded pipeline permitted <=1 repair call. Results (artifact-grounded): baseline = 149/150 (99.333\%); guarded = 150/150 (100.0\%); paired absolute difference = 0.006667 (\ensuremath{\approx}0.67 percentage points); exact McNemar two-sided p = 1.0; 95\% bootstrap percentile CI = [0.00, 0.02] (10,000 resamples, seed = 1729). Provider accounting: completed episodes = 300; model\_calls\_used = 151; repair-stage invocations = 1. Interpretation: on this preregistered synthetic corpus and under shared\_initial\_candidate semantics the guarded workflow produced a numerically negligible and statistically non-significant improvement over a near-ceiling baseline. We emphasize the methodological lesson that preregistered NetOps trials require baseline-calibration pilots to avoid ceiling effects that invalidate preregistered power assumptions. All execution and analysis artifacts are archived in the supplied execution bundle referenced by the execution manifest.
\end{abstract}

\section{Introduction}
Natural-language intent interfaces aim to reduce operator burden by letting humans express desired network changes as free text that is automatically translated into executable device configurations. Prior work documents that single-pass LLM outputs often contain syntactic errors, omissions, or semantic violations that can yield unsafe or unavailable states if applied directly [1][2][3]. A commonly proposed defense is a generate\ensuremath{\rightarrow}verify\ensuremath{\rightarrow}repair architecture in which an LLM produces a candidate, a deterministic verifier checks intent-level invariants or formalized constraints, and an automated repair (or human gate) corrects violations before execution [4][5]. Security studies of tool-augmented LLM agents further motivate deterministic gating because unchecked textual claims can become harmful actions in multi-tool workflows [6].

Research question. After an autonomous literature review and preregistration we posed a confined confirmatory question: does constrained prompting (structured templates with explicit semantic slots) combined with deterministic verifier-guided iterative feedback materially increase verifier-level semantic correctness of NL\ensuremath{\rightarrow}configuration translation versus an unconstrained baseline under a paired design? The trial (study\_cand\_2026\_b2\_v1) was preregistered and frozen before any model outputs were observed.

Contributions. (1) A fully preregistered, artifactized paired confirmatory trial measuring deterministic-verifier pass/fail on a programmatically generated synthetic corpus (N=150 paired intents). (2) Artifact-backed outcomes showing that, under shared\_initial\_candidate semantics and a one-repair budget, the guarded pipeline raised verifier pass rate from 149/150 to 150/150 (+0.67 pp), a numerically negligible and statistically non-significant improvement (exact McNemar p = 1.0; 95\% bootstrap CI = [0.00, 0.02]). (3) A methodological recommendation that preregistered NetOps trials include baseline-calibration pilots and explicitly specify generation semantics (shared\_initial\_candidate vs independent sampling) because semantics materially change the estimand and interpretation.

\section{Related Work}
Related literature spans intent-translation evaluations, verifier-guided repair systems, formal verifier tooling, and agentic/tool-augmented safety analyses. First, multiple recent studies evaluating NL\ensuremath{\rightarrow}configuration or intent-to-configuration pipelines report that translating underspecified operator language into verifier-sound artifacts is challenging in practice. Intent-driven translation work documents that free-text prompts frequently produce omissions or incorrect sequencing that violate required invariants, and published benchmark and evaluation studies report modest semantic-correctness rates on heterogeneous task banks [1]. Those evaluations emphasize that the mapping from natural intent to an exact required command sequence depends critically on prompt design, model competence, and the chosen correctness oracle.

Second, a growing body of work integrates deterministic verification into generation workflows and reports improved practical correctness, but gains depend strongly on baseline competence, verifier expressiveness, and repair budget. Systems that use policy-guided verifier feedback or verifier-in-the-loop fine-tuning (TerraFormer and related IaC efforts) demonstrate that verifier signals can be used to rerank or fine-tune candidate generations and thereby reduce syntactic and semantic errors for infrastructure-as-code artifacts [2]. Independent benchmark efforts focused specifically on LLM-driven configuration repair likewise show that generate\ensuremath{\rightarrow}verify\ensuremath{\rightarrow}repair pipelines can improve downstream validity, but reported effect sizes vary across datasets and experimental semantics (single-shot vs multi-sample generation, number of allowed repairs) and across whether the verifier feeds gradient-based updates or only deterministic textual feedback to a repair prompt [3]. These method papers collectively highlight two operational knobs that we treat explicitly in our design: the repair budget (how many additional generations/repairs are permitted) and the sampling semantics (whether different conditions draw independent first-pass samples or share an initial candidate). Both knobs materially affect the estimand and therefore the interpretable marginal benefit of verifier-guided repair.

Third, deterministic formal-verification tools provide concrete oracles for many network-level invariants but carry limitations that affect their role as evaluation endpoints. Fast symbolic/network-verification engines (e.g., NetKAT/KATch family) and related packet-reachability checkers can deterministically validate reachability, isolation, and many high-level policy constraints and have been proposed as suitability oracles for LLM-generated configurations [5]. Those tools are valuable as audit or gating layers because they return compact diagnostic counters and counterexamples that can drive automated repairs. At the same time, tool limitations---supported protocol models, topology scale, and abstraction mismatch with vendor-specific CLI sequences---mean that verifier coverage and mapping correctness become additional sources of error. Prior systems therefore combine careful mapping/unit-tests with verifier-coverage checks before committing to large-batch automated evaluation [5].

Fourth, security and systems analyses of agentic LLM pipelines warn that unguarded multi-tool agents can convert false textual claims into harmful actions and are susceptible to prompt-injection and claim-to-action exploit patterns [4]. Empirical attack studies show that deployed-style tool orchestration layers can be manipulated when the agent's decision pathways are not auditable or when verifier gating is absent; these findings motivate designs that make verifier verdicts explicit, auditable, and conservative in the presence of ambiguous intents.

How this study situates. Our trial operationalizes a narrowly constrained guarded architecture and emphasizes preregistration, paired inference, and deterministic-verifier endpoints. It differs from many prior reports in three design choices that prior work identifies as consequential: (a) we enforced shared\_initial\_candidate execution semantics so that baseline and guarded comparisons are computed on the identical sampled candidate (this isolates the marginal value of a repair applied to a fixed first-pass output); (b) we constrained the repair budget to at most one repair per intent, consistent with low-latency operational budgets considered in some IaC workflows; and (c) we used a deterministic verifier as the strict binary oracle for semantic correctness. These choices intentionally reduce the maximal observable marginal benefit of guarded repair under our estimand, explaining how previously reported larger improvements can coexist with our bounded result when different sampling or repair regimes are evaluated [2][3][5][6].

\section{Methodology}
Overview and preregistration. We preregistered study\_cand\_2026\_b2\_v1 and froze an analysis plan prior to any model calls. The design is paired-binary: each sampled intent yields one shared initial candidate used to evaluate both the baseline and guarded conditions (shared\_initial\_candidate semantics). The primary estimand was the paired\_success\_rate\_difference\_guarded\_minus\_baseline and the confirmatory test was an exact two-sided McNemar on the paired contingency table. The preregistration specified N=150 confirmatory intents, inclusion/exclusion rules, a one-repair budget for guarded episodes, a deterministic retry policy for provider timeouts (one retry), and an analysis plan that included a 10,000-resample paired bootstrap (bootstrap\_seed = 1729) to produce percentile CIs. These preregistered elements are archived; the registered preregistration SHA-256 is recorded in the execution manifest.

Task generation, intent families, and prechecks. Confirmatory tasks were produced deterministically by deterministic\_netops\_task\_generator\_v5. The generator emits intent payloads with: (i) an explicit natural-language intent string, (ii) an initial\_state and expected\_final\_state, (iii) an ordered required\_sequence of field-level operations (e.g., shutdown\ensuremath{\rightarrow}mtu\ensuremath{\rightarrow}restore), (iv) policy constraints (allowed\_touched\_fields, protected\_interfaces, final\_group\_constraints), and (v) a declared difficulty label. The confirmatory corpus was balanced across four intent families: safe sequence reachability/isolation, ACL-like access changes, VLAN/MTU sequence manipulations, and uplink switchover patterns. Topologies were restricted to verifier-capable small networks (4--32 nodes) to ensure deterministic validator coverage. Prior to model calls we ran deterministic unit-tests on the mapping/transformation code and a rule-based ambiguity detector; items flagged as ambiguous or failing parser/unit-tests were excluded per the preregistered exclusion rules. In this execution none of the confirmatory items were excluded.

Execution adapter, sampling semantics, and explicit baseline interpretation. The hosted\_netops\_gvr\_v1 adapter enforced the shared\_initial\_candidate semantics required by the preregistration: for each intent the adapter performed exactly one initial model generation (the shared candidate) and used that identical candidate when scoring both the baseline and guarded conditions. The guarded pipeline could issue at most one additional repair call only if the deterministic validator failed the shared initial candidate. Because of this architecture, the baseline success rate measures the validity of the single initial sampled candidate; the guarded vs baseline contrast therefore quantifies the marginal effect of a single allowed repair applied to the identical sampled candidate rather than the effect of independently re-sampling first-pass candidates per condition. We include this precise sampling-to-estimand mapping in Methods to eliminate interpretive ambiguity about whether different first-pass generations were drawn for each condition.

Model configuration and nondeterminism. Declared model: gpt-5-mini (provider record: openai\_responses). The archived execution/model\_configuration.json records declared\_model\_version = "gpt-5-mini" and explicitly records temperature = null (unset). Importantly, the adapter and preregistration did not enforce a numeric temperature value or a fixed random seed; therefore model outputs may be nondeterministic across independent API sessions. We state this here in Methods and repeat it in the Disclosure Statement to comply with preregistration/runtime-parameter reporting requirements.

Repair budget, session isolation, and provider auditing. The guarded repair policy allowed at most one repair invocation per intent and the adapter maintained strict session isolation: each model call ran in a fresh API session (no cross-call state reuse). Provider auditing tracked hosted\_model\_call\_event\_count and cache statistics; the deterministic reconciliation artifact reports completed\_hosted\_model\_call\_event\_count = 151, cache\_miss\_count = 151, and stage\_counts = \{shared\_initial\_generation: 150, repair: 1\}. The adapter implementation enforces the preregistered retry rule (one retry after transient provider failure). All provider-call metadata and per-call runtime parameters were archived for audit.

Deterministic validator, scoring rule, and diagnostic trace structure. We used deterministic\_netops\_validator\_v5 (validator\_version = 5.0) as a deterministic oracle for the primary binary endpoint. The validator implements a full\_intent\_constraint\_validation rule that: (a) parses the LLM-produced candidate into normalized command sequences, (b) computes parsed\_command\_count and compares against required\_command\_count derived from the task's required\_sequence, (c) enforces workflow and group constraints (e.g., make-before-break, must\_be\_down\_for\_fields), and (d) returns a boolean valid flag plus structured diagnostics (violation\_codes, observed\_touched\_field\_count, normalized\_configuration, final\_state snapshot). Scoring for the confirmatory endpoint used the validator's boolean pass/fail; the richer diagnostics were used for failure-mode analysis and are archived as validator traces for each episode.

Inclusion/exclusion, contamination controls, and pre-execution unit-tests. The preregistered inclusion rules required that an intent pass the deterministic ambiguity detector and that mapping code successfully produce verifier inputs. We executed deterministic unit-tests for mapping code and verifier-input transformations before batch evaluation; these tests passed. Contamination checks (session-leak detection, duplicate-response entropy checks) were run; contamination\_summary.csv recorded zero flags. Deterministic reconciliation confirms episode accounting consistency and reports no provider-call failures.

Statistical analysis posture (confirmatory vs exploratory). The methodology adheres to the preregistered multiplicity and missingness rules: the primary confirmatory contrast is guarded vs baseline at \ensuremath{\alpha}=0.05 (exact McNemar two-sided) using complete\_pair\_primary (exclude incomplete pairs). Secondary contrasts were preregistered (constrained vs baseline, guarded vs constrained) with Bonferroni correction; in practice shared\_initial\_candidate semantics and adapter limits constrained the realized secondary contrasts. Analysis artifacts include paired\_contingency\_table.csv, condition\_summary.csv, and a paired bootstrap used to compute the 95\% CI (bootstrap\_resamples = 10000, bootstrap\_seed = 1729). All analysis code, seeds, and artifacts are archived for reproducibility.

Reproducibility artifacts and archival pointers. The full execution manifest, raw model-call logs, per-call runtime metadata, task-generator seeds, verifier inputs/verdicts, and analysis outputs are archived in the execution bundle referenced by the execution manifest. Deterministic reconciliation and provider accounting artifacts record model\_calls\_used = 151 and stage-level counts; these artifacts support replication of the paired accounting and the precise sampling semantics used in this trial.

\section{Results}
Execution and provider audit. The run completed as planned: completed\_episode\_count = 300 and complete\_pair\_count = 150. Provider-call accounting: model\_calls\_used = 151; completed\_hosted\_model\_call\_event\_count = 151; failed\_hosted\_model\_call\_event\_count = 0; cache\_miss\_count = 151. Stage-level counts: shared\_initial\_generation = 150 and repair = 1. No contamination flags or pre-execution unit-test failures occurred.

Primary numerical outcomes and contingency counts. The paired contingency counts (guarded \ensuremath{\times} baseline) were: n\_11 = 149 (both-pass), n\_10 = 1 (guarded-only-pass), n\_01 = 0 (baseline-only-pass), n\_00 = 0 (both-fail). Per-condition success rates: baseline\_success\_count = 149/150 = 0.9933333333333333; guarded\_success\_count = 150/150 = 1.0. Estimated paired absolute difference (guarded - baseline) = 0.00666666666666671 (\ensuremath{\approx}0.67 percentage points). Exact McNemar two-sided p = 1.0. Paired bootstrap percentile 95\% CI = [0.00, 0.02] (10,000 resamples; bootstrap\_seed = 1729). These figures are reported verbatim from the archived analysis artifacts (analysis/results.json, analysis/paired\_contingency\_table.csv, analysis/condition\_summary.csv).

Repair diagnostics and revision iterations. Only one guarded repair invocation occurred (1/150), consistent with the near-ceiling initial success rate. That single repair converted one initially invalid candidate to a valid configuration under the deterministic verifier. Median revision iterations for guarded episodes = 0 (IQR = 0). Validator traces include validation\_before and validation\_after states, normalized\_configuration strings, parsed\_command\_count, required\_command\_count, violation\_codes, and a valid flag for each episode; representative traces for tasks 000001--000006 are archived and demonstrate varied difficulty patterns while confirming validator behavior.

Representative artifact-grounded example. Example pair (task-000001, generator\_seed = 1) intent: "Migrate edge1 to VLAN 30 and MTU 1600. For safety, shut edge1 down before changing either VLAN or MTU, then restore it to admin up. Do not modify any other interface." Shared-initial response and both condition responses returned the expected four-command sequence. Validator diagnostics for task-000001: parsed\_command\_count = 4, required\_command\_count = 4, observed\_touched\_field\_count = 3, and valid = true. The shared-initial candidate SHA-256 and normalized\_configuration strings are archived in the execution bundle as part of the scorer traces.

Sensitivity and missingness. The preregistered sensitivity analysis treating missing guarded outputs as failures was not invoked because failed\_hosted\_model\_call\_event\_count = 0 and unscorable episodes = 0. No exclusions or parse failures occurred and contamination\_summary.csv is empty.

Compact evidence tables. Table 1 (paired contingency) and Table 2 (execution audit summary) summarize the confirmatory counts and provider accounting; both tables are artifact-grounded and constructed from analysis/paired\_contingency\_table.csv and analysis/condition\_summary.csv.

Interpretation of numeric outcome. On this synthetic corpus and under shared\_initial\_candidate semantics the guarded workflow produced a numerically negligible absolute improvement (+0.67 pp) from an already near-ceiling baseline. The exact McNemar p = 1.0 and the bootstrap CI containing zero indicate no statistically supported confirmatory benefit under the preregistered estimand. Importantly, the observed baseline success rate (\ensuremath{\approx}99.33\%) violated the preregistered power assumption (target detectable effect = 15 pp), producing a ceiling effect that invalidated the originally planned sensitivity.

\section{Discussion}
Ceiling effect and implications for preregistration. The preregistered power calculation assumed a substantially lower baseline; observed baseline\_success\_rate = 0.9933 left almost no headroom for the guarded repair to show a large absolute improvement. This ceiling effect invalidated the preregistered detectable-effect assumptions (target = 15 percentage points) and reduced the trial's ability to demonstrate benefit. We therefore recommend that preregistered NetOps trials include a short baseline-calibration pilot (e.g., 20--40 tasks sampled from the planned generator and prompt style) to estimate baseline competence and adjust N or task difficulty before committing to confirmatory sampling. The archived artifacts contain the deterministic task generator seeds and representative pilot-able tasks for such calibration.

Semantics-driven estimand differences and practitioner implications. Under shared\_initial\_candidate semantics the baseline measures the single initial sampled candidate while the guarded condition quantifies the marginal value of at most one repair applied to that same candidate. Independent-sampling designs or multiple-repair budgets yield different operational estimands---expected performance averaged over independent draws or repair-driven robustness gains---and can produce larger observable improvements when baseline sampling variance is larger. Practitioners should choose sampling semantics that align with operational constraints (e.g., whether the deployed system will re-sample across alternatives or apply repairs to a fixed first-pass candidate).

Operational affordances beyond the binary pass/fail metric. Deterministic guarded workflows provide operational benefits not captured by a single binary endpoint: auditable normalized configurations for operator review, concise diagnostic traces (parsed\_command\_count, violation\_codes) that accelerate triage, and deterministic verifier verdicts suitable for policy gating. Even when marginal verifier-pass improvements are small on constrained corpora, these operational affordances can justify guarded architectures in practice; the archived validator traces demonstrate the concrete diagnostic strings that would be available to operators.

Threats to validity. Internal validity: adapter-enforced call budgets and provider audit mitigate cross-condition leakage; the shared\_initial\_candidate semantics intentionally reduce sampling variance but do not capture per-condition randomness. Construct validity: deterministic validator pass/fail is a proxy for semantic correctness and depends on mapping code and verifier coverage; mapping risks were mitigated with unit-tests and archived artifacts but residual risk remains. External validity: experiments used a single declared model (gpt-5-mini), synthetic tasks, and small topologies (4--32 nodes); claims are limited to the archived corpus. Conclusion validity: preregistered power assumptions were invalidated by the high baseline, constraining confirmatory inference. All such limitations are reported and linked to archived artifacts and counts.

Recommendations and future experimental designs. Based on the artifact-grounded outcomes we recommend: (1) baseline-calibration pilots to select N and difficulty; (2) explicit preregistration of generation semantics (shared\_initial\_candidate vs independent) and repair budgets; (3) repair-budget arms (multiple allowed repairs) to estimate diminishing returns; (4) model-diversity arms to assess generality across model families and sampling parameters; and (5) incremental scaling of verifier coverage and topology size with unit-tested mapping code before batch execution. The execution manifest and representative task set provide seeds and code to implement these next-step designs reproducibly.

Reproducibility and artifact availability. All code, task-generator seeds, model-configuration metadata, raw model-call logs, verifier inputs/verdicts, and analysis artifacts (contingency tables, bootstrap parameters, and provider-call audits) are archived in the execution bundle referenced by the execution manifest. Key numeric analysis parameters are recorded (bootstrap\_seed = 1729; bootstrap\_resamples = 10000), and all analysis CSVs and validator traces cited above are available for independent verification. The archived initial\_master\_prompt reference and preregistration record are included as immutable artifacts.

\section{Limitations}
\begin{itemize}
\item Single-model constraint: experiments used only gpt-5-mini; results do not generalize across models or sizes.
\item Synthetic corpus and scale: tasks were programmatically generated and limited to verifier-capable toy-to-small topologies (4--32 nodes); external validity to production WANs and heterogeneous vendor CLIs is untested.
\item Shared\_initial\_candidate semantics and execution contract limited repair budget to at most one guarded repair call per intent; different generation semantics or multiple-repair budgets may yield different outcomes.
\item Ceiling effect and preregistration mismatch: baseline correctness was far higher than the pilot assumptions used for power calculations (planned detectable effect = 15 percentage points), reducing the trial's ability to demonstrate benefit and illustrating sensitivity to baseline calibration.
\item Verifier coverage and specification translation: the study measures deterministic validator pass/fail on the preregistered invariants but does not claim safety guarantees beyond the deterministic validator's modeled properties.
\item Security and adversarial robustness: the experiment did not evaluate adversarial prompt-injection or adversary-driven intents; separate targeted studies are required to assess claim-to-action vulnerabilities in agentic NetOps workflows.
\end{itemize}

\section{Conclusion}
We executed a fully preregistered paired confirmatory trial comparing baseline free-text prompting with constrained-template prompting plus a single verifier-guided repair under shared\_initial\_candidate semantics. On a deterministic synthetic corpus with gpt-5-mini and deterministic\_netops\_validator\_v5, the guarded pipeline increased verifier pass rate from 149/150 to 150/150 (\ensuremath{\approx}0.67 pp), a numerically tiny and statistically non-significant difference (exact McNemar p = 1.0; 95\% bootstrap CI = [0.00, 0.02]). The principal scientific lesson is methodological: preregistered NetOps trials should include baseline calibration to avoid ceiling effects and preserve statistical power. Technically, our artifacts bound the marginal benefit of a one-repair guarded pipeline under shared\_initial\_candidate semantics; evaluating independent-sampling semantics, larger repair budgets, model diversity, and production-scale topologies remains important future work.

\section*{Disclosure Statement}
This study was executed autonomously after an autonomy lock under the archived master prompt (provenance/master\_prompt.txt, SHA-256 = 1872df1e\allowbreak{}1805d2d9\allowbreak{}6940456c\allowbreak{}a016bd66\allowbreak{}5d1d5196\allowbreak{}add77f5a\allowbreak{}cdf1582b\allowbreak{}b39b15ba). Preregistration identifier: study\_cand\_2026\_b2\_v1 (preregistration was frozen prior to observing outcomes and the preregistration record is archived). Declared model/tooling: declared model = gpt-5-mini (provider record: openai\_responses); execution adapter family = hosted\_netops\_gvr\_v1; deterministic validator = deterministic\_netops\_validator\_v5 (validator\_version = 5.0). Runtime sampling semantics (explicit): the archived model\_configuration.json records temperature = null (unset); because no numeric temperature or fixed random-seed was enforced, model outputs may be nondeterministic across independent API sessions. Autonomy and human role: a human Research Director selected the topic family and constrained the initial master prompt prior to the autonomy lock; after the lock the autonomous system performed literature review, study selection, preregistration, code generation, experiment execution, analysis, manuscript drafting, autonomous peer review, and revision. No human manually edited technical content, analyses, figures, captions, or references after the autonomy lock. Key execution facts reported in the paper (artifact-grounded): completed episodes = 300; complete paired tasks = 150; model\_calls\_used = 151; repair-stage invocations = 1. All runtime sampling metadata, provider-call traces, verifier inputs/verdicts, and analysis artifacts are archived in the execution bundle referenced by the execution manifest.

\begin{thebibliography}{99}
\bibitem{ref1} https://doi.org/10.1002/nem.2313
\bibitem{ref2} https://doi.org/10.1145/3786583.3786898
\bibitem{ref3} https://arxiv.org/abs/2604.22513
\bibitem{ref4} https://doi.org/10.1145/3605764.3623985
\bibitem{ref5} https://doi.org/10.1145/3656454
\bibitem{ref6} https://doi.org/10.1002/widm.70111
\end{thebibliography}

\end{document}
